\documentclass[12pt]{article}
\usepackage[margin=1in]{geometry}
\usepackage{enumitem}
\usepackage{titlesec}
\usepackage{parskip}
\usepackage{graphicx}
\usepackage{hyperref}
\usepackage{fontenc}

\title{\textbf{Weber 1930 RG Notes}\\
\large Weber, Max. \textit{The Protestant Ethic and the Spirit of Capitalism}. Trans. Talcott Parsons.\\
London: George Allen \& Unwin, 1930 [1904--1905].}
\author{}
\date{}

\begin{document}

\maketitle

\section{Main Idea}



%--------------------------------------------------
\section{General Notes}
%--------------------------------------------------

\begin{itemize}[leftmargin=*, itemsep=4pt]
    \item Some similarity between Marx and Weber, in that both were seeking to explain modern economic development and consider the conditions it both required and created -- but, Weber took specific pieces of Marxist ideology while explicitly rejecting its total implementation/acquisition. Weber was particularly critical of some of the Marxist followers of this time and especially their revolutionary sentiments; instead sought a much more peaceful transition and to explain the organic development of the modern economy.
    \item Though desires for wealth acquisition are not novel, the contemporary West has become associated with ``the rational organisation of formally free labour.'' -- its ``routinised, calculated administration within continuously functioning enterprises.'' This implies both a disciplined labour force and the regularised investment of capital. (Quote from interpreter but not from Weber)
    \item At this time, economic development and wealth accumulation seemed to be highly correlated with the presence of Protestant communities: such societies were achieving much greater economic development than their non-Protestant counterparts; in particular, Catholics.
    \item Weber claims that the underlying connection between Protestantism and a developed work ethic -- particularly through the ``virtuous'' reinvestment of capital and pursuit of meaningful labor -- were particularly key components in the development of modern economies (and in the heterogeneity of such development). Everyday men sought to pursue meaningful work, which in turn led to the development of massive amounts of capital and development
\end{itemize}


\section{Important Specific Factual Claims}

\begin{itemize}[leftmargin=*, itemsep=4pt]
    \item In a context of predestination post-Calvin, people sought to prove externally that they were to be saved, and performed worldly ``good works'' as a demonstration of surety. Subsequently, fulfilling one's calling was no longer a means by which salvation could be achieved, but a signal to others that the individual was to be saved.
\end{itemize}

\section{Connections}

\begin{itemize}
    \item Lipset 1959: ``For the purpose of this paper, democracy (in a complex society) is defined as a political system which supplies regular constitutional opportunities for changing the governing officials. It is a social mechanism for the resolution of the problem of societal decision-making among conflicting interest groups which permits the largest possible part of the population to influence these decisions through their ability to choose among alternative contenders for political office. In large measure abstracted from the work of Joseph Schumpeter and Max Weber, this definition implies a number of specific conditions: (a) a ``political formula,'' a system of beliefs, legitimizing the democratic system and specifying the institutions -- parties, a free press, and so forth -- which are legitimized, i.e. accepted as proper by all; (b) one set of political leaders in office; and (c) one or more sets of leaders, out of office, who act as a legitimate opposition attempting to gain office.'' Talks quite a bit more about Weber as well
\end{itemize}

\section{My Critiques}

\begin{itemize}
    \item
\end{itemize}


\end{document}
